\documentclass[12pt,a4paper]{article}
\usepackage[margin=1in]{geometry}
\usepackage{amsmath,amssymb,amsthm}
\usepackage[hidelinks]{hyperref}
\usepackage{booktabs}
\usepackage{array}
\newcolumntype{P}[1]{>{\raggedright\arraybackslash}p{#1}}

\title{The Imprint Chain:\\
A Roadmap and Epistemic Grading of a Research Programme}

\author{
  Sidong Liu, PhD \\
  iBioStratix Ltd \\
  \texttt{sidongliu@hotmail.com}
}

\date{July 2026}

\begin{document}
\emergencystretch=2em
\raggedbottom
\hbadness=5000
\vbadness=5000

\maketitle

\begin{abstract}
This paper presents a roadmap rather than a new theorem. Its purpose is to display and grade a chain that runs from algebraic emergence through scale-stable imprinting and accessible world-formation to self-reading mind and the open question of finite self-description. The central claim is that the chain is useful as a research architecture, but its segments do not all have the same epistemic status.

The paper's two practical contributions are an explicit epistemic grading map and a set of reading paths for external readers who do not want to traverse the full body of work in chronological order.

\medskip

\noindent\textbf{Keywords}: emergence, imprint, self-reading, roadmap, programme-level bridge, epistemic grading, consciousness
\end{abstract}

\section{Introduction: Why a Roadmap Is Needed}

The body of work linking the Genesis Trilogy, the Consciousness Series, and \emph{The World as Imprint} now contains enough internal structure that external readers face an avoidable difficulty. The difficulty is not primarily one of disagreement. It is one of entry. Which claims are supposed to be read as formal results? Which are only programme-level bridges? Which are interface proposals between technical and philosophical languages? Which are structural resonances that orient interpretation but do not function as evidence?

This paper is written to answer those questions directly. It does not add a new mechanism to the chain from emergence to self-reading mind. It displays the chain, grades its segments, and marks where different kinds of claim meet.

The need for such a map became especially clear once \emph{The World as
Imprint} introduced a compressed philosophical vocabulary capable of
speaking across several registers \cite{WorldImprint}. That vocabulary
now treats imprint-carrying as decoder-relative and possibly null, while
self-reading remains a separate programme-level threshold. The bridge
papers clarify local interpretive links; they do not inherit a theorem
from the imprint heuristic.

The paper's governing sentence is therefore the following:

\begin{quote}
This paper does not prove the chain; it displays it, grades it, and marks the places where theorem-level results, programme-level bridges, interface claims, and structural resonances meet.
\end{quote}

That sentence expresses both the ambition and the restraint of the paper. It is ambitious because it attempts to present one coherent architecture from source-side emergence to finite self-description. It is restrained because it refuses to blur the difference between a theorem and a philosophical bridge.

The grading discipline is not arbitrary. It resonates with structural-realist programmes in philosophy of science, which separate structural content from extra-structural commitments \cite{Ladyman2007,French2014}. It also sits naturally beside Batterman's emphasis on asymptotic and cross-scale explanation, where explanatory power often outruns one-to-one microstructural reconstruction \cite{Batterman2002}. This paper extends that orientation by adding an explicit category of interface claims between mature vocabularies.

\section{The Full Imprint Chain}

The chain can be written in its shortest usable form:

\begin{quote}
source algebra $\to$ coarse-graining $\to$ scale-stable imprints $\to$
imprint-field $\to$ integration workspace \(\mathcal W_{\mathrm{int}}\)
$\to$ self-model $\to$ self-reading threshold $\dashrightarrow$
open finite-self-description question.
\end{quote}

This sequence is not meant as a single deduction. It is a structured path across multiple already developed works. Each arrow asks a different question:

\begin{center}
\begin{tabular}{P{0.34\linewidth}P{0.56\linewidth}}
\toprule
\textbf{Segment} & \textbf{Question} \\
\midrule
source algebra $\to$ coarse-graining & How does lower-level structure become effectively describable at all? \\
coarse-graining $\to$ scale-stable imprints & Which marks survive scale change as readable residues? \\
scale-stable imprints $\to$ imprint-field & How should emergence be read once its surviving residues are treated as trace-carrying? \\
imprint-field $\to\mathcal W_{\mathrm{int}}\to$ self-model & Under what additional cognitive interface do readable traces become an integrated world supporting self-location? \\
self-model $\to$ self-reading threshold & When does the model become reflexively available as itself imprint-bearing? \\
self-reading threshold $\dashrightarrow$ finite self-description horizon & Can a physical horizon be derived from a specified observation or reconstruction model? \\
\bottomrule
\end{tabular}
\end{center}

The paper's first task is simply to insist that these are not all the same kind of question. Confusion begins whenever they are treated as if they belonged to one epistemic register.

\section{Segment I: Algebraic Emergence}

The first segment concerns the move from source-side structure to coarse-grained observability. This is largely the territory of the Genesis Trilogy and of Foundation~I \cite{HAFFA,HAFFD,FoundationI}. The relevant lesson is not that every source-side detail survives. It is that emergence is already filtered by accessibility, projection, or effective selection.

This segment should be read as a mixture of formal and conditional work. In the stronger Genesis papers, coarse-grained observable structure is treated in explicit algebraic terms. In Foundation~I, the passage from source datum to accessibility remains conditional in precise ways. The point for the present map is that the chain does not begin with already-given minds or worlds. It begins with the problem of how any effective descriptive regime becomes available at all.

\section{Segment II: Scale-Stable Imprint Filtering}

The second segment is the task of the first bridge paper, \emph{Renormalization as Imprint Filtering}. Its role is to explain what stability means in the imprint vocabulary once scale change is taken seriously. The claim here is programme-level rather than theorem-level: renormalization language can be read as answering which readable residues survive coarse-graining.

This segment is important because without it the imprint vocabulary risks remaining too intuitive. ``Stable'' would then sound like a decorative adjective rather than a physically constrained notion. The bridge paper sharpens it by relating stability to relevance, universality, fixed-point behavior, and symmetry-breaking residue.

\section{Segment III: Decoder-Relative Imprint-Field}

The third segment is the conceptual core of \emph{The World as Imprint} \cite{WorldImprint}. Once scale-stable residues are in view, the emergent world can be investigated as a candidate imprint-field: not a hidden substrate, but a decoder-relative domain whose imprint value may be nonzero, small, or null.

This segment is explicitly programme-level and philosophical. It is not a theorem, and it does not pretend to be one. Its contribution is to give a disciplined vocabulary for a structural fact already distributed across the more formal papers: emergence is not blank outcome but marked organization.

\section{Segment IV: Stable Hull Integration}

The fourth segment is the task of the second bridge paper,
\emph{Stable Hulls as Imprint Integration}. Its live object is the
internal cognitive workspace \(\mathcal W_{\mathrm{int}}\). This is
distinct from Foundation~I's source-algebraic
\(\mathcal H_{\mathrm{st}}\) and Foundation~II's external shared sector
\(\mathcal H_{12}\) \cite{FoundationII,BrainAnchor}. No Foundation
construction is imported as the cognitive workspace.

This segment is an interface claim. It proposes a middle cognitive
architecture between decoder-relative imprint and reflexive recognition;
it is not a theorem inherited from Foundation~II or from the imprint
heuristic.

The self-model literature is developed more directly in \emph{Stable Hulls as Imprint Integration} \S6; this roadmap only locates that engagement.

\section{Segment V: The Self-Reading Threshold}

The fifth segment returns to \emph{The World as Imprint}. There the central distinction is between an imprint-carrier and a self-reading system \cite{WorldImprint}. The self-reading threshold is crossed when a system can reflexively read its own world-model, or itself, as imprint-bearing.

This segment should again be read as programme-level rather than theorem-level. It provides a marker matrix and a conceptual threshold, not a full constitutive theory of first-person consciousness. Its value lies in holding apart several things often confused: structure, responsiveness, self-modeling, and reflexive self-reading.

\section{Segment VI: Self-Description Status}

The final segment asks whether organized self-reading has a physically
derived finite horizon. Consciousness Paper~III formerly answered this
with a local minimax drift floor and a physical control order
\(k(n)\), but both claims have been retracted \cite{PaperIII}.

The only surviving formal result in that segment is the conditional
arithmetic lemma
\[
\log I_k=-\alpha k\log k+O(k)
\quad\Longrightarrow\quad
K(n)=O\!\left(\frac{\log n}{\log\log n}\right)
\]
under a stipulated positive sequence and stipulated budget. It proves
neither a physical Fisher spectrum nor a self-model budget. A finite
self-description horizon is therefore conditional/open, not a proved
endpoint of the chain.

\section{Epistemic Grading Map}

The chain becomes clearer once its internal grades are stated explicitly.
The status attaches to an arrow, not automatically to either endpoint.

\begin{center}
\begin{tabular}{P{0.18\linewidth}P{0.72\linewidth}}
\toprule
\textbf{Arrow status} & \textbf{Meaning} \\
\midrule
\textbf{THM} & Verified theorem: the arrow itself is proved, logically complete, and has no open dependency. \\
\textbf{COND} & Conditional theorem: the arrow is proved after all named assumptions or external inputs are stated. \\
\textbf{MODEL} & Model check: the connection is verified only in a declared finite or concrete model and is not extrapolated. \\
\textbf{IAX} & Interface axiom: the target explicitly supplies the cross-paper connection as an input rather than deriving it. \\
\textbf{DEF} & Definition-unpacking: the target use follows by expanding a shared definition and adds no substantive implication. \\
\textbf{HEUR} & Heuristic bridge: the connection is a disciplined analogy, interpretation, or research guide, not an inference. \\
\textbf{META} & Metaphor or programme-level orientation with no evidential or mathematical role. \\
\textbf{INVALID-OPEN} & The asserted arrow is withdrawn, contradicted, or missing a required construction; no current implication is available. \\
\bottomrule
\end{tabular}
\end{center}

The formal results at individual nodes remain important, but they do not
make a cross-paper arrow formal by inheritance. The following two
registers therefore separate the seven narrative arrows of the imprint
chain from the seventeen dependency-matrix edges audited across the
programme.

\subsection{Roadmap Narrative-Arrow Register}

\subsubsection*{P1. Source algebra to effective coarse-graining}
\noindent\textbf{Arrow:} Genesis/Foundation source structure
\(\to\) effective coarse-grained description.\par
\noindent\textbf{Source Claim ID:} \texttt{HAFF-A-T2.6},
\texttt{HAFF-A-T3.4}, and \texttt{FI-hull}.\par
\noindent\textbf{Source current status:} \texttt{HAFF-A-T2.6} is
\texttt{DOWNGRADE}; \texttt{HAFF-A-T3.4} is a conditional model result;
\texttt{FI-hull} is \texttt{OK-cond} under its invariant-class,
modular-invariance, and expectation hypotheses.\par
\noindent\textbf{Target usage description:} Opens the roadmap with the
question of how an effective descriptive regime is selected.\par
\noindent\textbf{Arrow status:} \textbf{HEUR}.\par
\noindent\textbf{Failure effect if mis-graded:} A composite organizational
transition would be mistaken for a universal derivation of accessibility
or geometry.\par
\noindent\textbf{Notes:} The surviving node-level results have different
domains and conditions; no single theorem identifies them with this
compressed arrow.\par

\subsubsection*{P2. Coarse-graining to scale-stable imprint}
\noindent\textbf{Arrow:} Ordinary RG flow facts \(\to\) scale-stable
imprint vocabulary.\par
\noindent\textbf{Source Claim ID:} \texttt{RG-map}.\par
\noindent\textbf{Source current status:} \texttt{HEUR/DOWNGRADE};
relevant directions grow, irrelevant directions remain as corrections,
and universality is many-to-one.\par
\noindent\textbf{Target usage description:} Rereads selected RG residues
as candidate decoder-readable marks.\par
\noindent\textbf{Arrow status:} \textbf{HEUR}.\par
\noindent\textbf{Failure effect if mis-graded:} The bridge would falsely
claim that RG proves imprint preservation or microscopic reconstruction.\par
\noindent\textbf{Notes:} No theorem is inherited from \texttt{WI-princ},
and the decoder-relative value may be null.\par

\subsubsection*{P3. Scale-stable imprint to imprint-field}
\noindent\textbf{Arrow:} Candidate scale-stable residue \(\to\)
decoder-relative imprint-field.\par
\noindent\textbf{Source Claim ID:} \texttt{WI-princ}.\par
\noindent\textbf{Source current status:} \texttt{HEUR/DOWNGRADE}; the
operational target requires a source family, decoder class, score, and a
genuine zero case.\par
\noindent\textbf{Target usage description:} Supplies a programme-level
vocabulary for asking which emergent structures carry readable source
contrasts.\par
\noindent\textbf{Arrow status:} \textbf{HEUR}.\par
\noindent\textbf{Failure effect if mis-graded:} Decoder-relative and
possibly null carrying would be inflated into universal nonzero imprint.\par
\noindent\textbf{Notes:} The target is
\(\operatorname{Imp}_{\mathcal D}(\Theta;\mathcal S_\mu)\), a
definition-to-be-discharged rather than an established invariant.\par

\subsubsection*{P4. Imprint-field to internal integration workspace}
\noindent\textbf{Arrow:} Decoder-relative imprint vocabulary plus access
assumptions \(\to\mathcal W_{\mathrm{int}}\).\par
\noindent\textbf{Source Claim ID:} \texttt{WI-princ},
\texttt{SH-central}, and \texttt{BRAIN-P6.1}.\par
\noindent\textbf{Source current status:} \texttt{WI-princ} and
\texttt{SH-central} are downgraded heuristics; the Brain construction now
has explicit conditional hypotheses and a finite model witness, while
the general private-complement package remains open.\par
\noindent\textbf{Target usage description:} Proposes a cognitive
integration regime for readable traces inside one accessible algebra.\par
\noindent\textbf{Arrow status:} \textbf{HEUR}.\par
\noindent\textbf{Failure effect if mis-graded:} The cognitive workspace
would be falsely derived from imprint vocabulary or silently identified
with \(\mathcal H_{\mathrm{st}}\) or \(\mathcal H_{12}\).\par
\noindent\textbf{Notes:} The Brain witness checks one
Pauli-qubit model; it does not turn this interpretive bridge into a
cross-series theorem.\par

\subsubsection*{P5. Workspace to self-model}
\noindent\textbf{Arrow:} Internal workspace
\(\mathcal W_{\mathrm{int}}\to\) reportable self-model.\par
\noindent\textbf{Source Claim ID:} \texttt{BRAIN-P6.1}.\par
\noindent\textbf{Source current status:} Contract status
\texttt{DOWNGRADE}; the revised source separates a conditional
fixed-point construction, a finite Type~I witness, and a biological
interface reading.\par
\noindent\textbf{Target usage description:} Interprets the exact common
workspace as a candidate structure supporting report and self-location.\par
\noindent\textbf{Arrow status:} \textbf{HEUR}.\par
\noindent\textbf{Failure effect if mis-graded:} A finite algebraic model
would be presented as a neuroscience theorem or a proof of subjectivity.\par
\noindent\textbf{Notes:} The model check establishes non-vacuity of the
formal package only; the biological interpretation remains interface-level.\par

\subsubsection*{P6. Self-model to self-reading threshold}
\noindent\textbf{Arrow:} Reportable self-model \(\to\) self-reading
threshold.\par
\noindent\textbf{Source Claim ID:} \texttt{WI-princ}.\par
\noindent\textbf{Source current status:} \texttt{HEUR/DOWNGRADE}.\par
\noindent\textbf{Target usage description:} Uses a marker matrix to
distinguish ordinary self-modeling from reflexive recognition of
constraint-shaped participation.\par
\noindent\textbf{Arrow status:} \textbf{HEUR}.\par
\noindent\textbf{Failure effect if mis-graded:} A programme-level marker
would become a necessary or sufficient constitutive criterion for
consciousness.\par
\noindent\textbf{Notes:} The threshold organizes questions; it is not a
validated psychometric or phenomenological boundary.\par

\subsubsection*{P7. Self-reading to a finite physical horizon}
\noindent\textbf{Arrow:} Self-reading threshold
\(\dashrightarrow\) finite self-description horizon.\par
\noindent\textbf{Source Claim ID:} \texttt{CIII-T3.3},
\texttt{CIII-kn}; surviving arithmetic node \texttt{CIII-asym}.\par
\noindent\textbf{Source current status:} \texttt{CIII-T3.3} and
\texttt{CIII-kn} are \texttt{RETRACT}; \texttt{CIII-asym} is
\texttt{OK-cond} only as an abstract arithmetic implication.\par
\noindent\textbf{Target usage description:} Would supply a physical
endpoint for the imprint chain.\par
\noindent\textbf{Arrow status:} \textbf{INVALID-OPEN}.\par
\noindent\textbf{Failure effect if mis-graded:} The roadmap would restore
the refuted positive drift floor and physical \(k(n)\) theorem.\par
\noindent\textbf{Notes:} The arithmetic lemma establishes neither a
physical Fisher spectrum nor a self-model budget.\par
\noindent\textbf{Repair path:} Construct a typed observation model that
excludes lossless-orbit codes, prove the estimator-risk/reconstruction-loss
link, and independently derive normalized physical coefficient
directions and a resource budget.\par

\subsection{Programme Dependency-Edge Register}

\subsubsection*{E1. Former global accessible algebra to its consumers}
\noindent\textbf{Arrow:} \texttt{HAFF-E-D3.10}
\(\to\) HAFF B/D/F, T-DOME I, and Foundation I uses.\par
\noindent\textbf{Source Claim ID:} \texttt{HAFF-E-D3.10}.\par
\noindent\textbf{Source current status:} \texttt{DEF/RETRACT}; the old
\(\mathcal A_c\) definition lacks identity and closure and conflicts with
its own criteria.\par
\noindent\textbf{Target usage description:} Formerly supplied one global
accessible algebra to several differently typed consumers.\par
\noindent\textbf{Arrow status:} \textbf{INVALID-OPEN}.\par
\noindent\textbf{Failure effect if mis-graded:} Probe effects, pointer
records, fixed points, selected geometric algebras, and stable hulls would
again be conflated.\par
\noindent\textbf{Notes:} Current consumers have been rebuilt with
separate typed objects.\par
\noindent\textbf{Repair path:} No repair to the global-substitution edge
is warranted; retain the consumer-specific reconstructions unless a new
typed universal construction is proved.\par

\subsubsection*{E2. Withdrawn uniqueness to multiplicity narrative}
\noindent\textbf{Arrow:} \texttt{HAFF-D-T4.6}
\(\to\) HAFF A/B multiplicity narrative.\par
\noindent\textbf{Source Claim ID:} \texttt{HAFF-D-T4.6}.\par
\noindent\textbf{Source current status:} \texttt{THM/RETRACT}; the
ambient modular flow is undefined in the stated pure-state setting.\par
\noindent\textbf{Target usage description:} Retains local
underdetermination and candidate multiplicity as motivation only.\par
\noindent\textbf{Arrow status:} \textbf{HEUR}.\par
\noindent\textbf{Failure effect if mis-graded:} The withdrawn uniqueness
theorem would be used to manufacture a uniqueness-versus-multiplicity
reconciliation.\par
\noindent\textbf{Notes:} The local multiplicity observations do not
depend on the failed uniqueness result.\par

\subsubsection*{E3. Final-cut diagnosis to no-view closure}
\noindent\textbf{Arrow:} \texttt{HAFF-G-L3.1}
\(\to\) \texttt{QR-NVN-P1} closure principle.\par
\noindent\textbf{Source Claim ID:} \texttt{HAFF-G-L3.1}.\par
\noindent\textbf{Source current status:} \texttt{META/RETRACT}; no
fixed-point exclusion theorem is proved.\par
\noindent\textbf{Target usage description:} Supplies programme-level
motivation for a separately stated no-view diagnosis.\par
\noindent\textbf{Arrow status:} \textbf{HEUR}.\par
\noindent\textbf{Failure effect if mis-graded:} Two downgraded
meta-closure claims would appear to prove each other.\par
\noindent\textbf{Notes:} Neither endpoint formally supports the other.\par

\subsubsection*{E4. Clifford model choices to T-DOME II}
\noindent\textbf{Arrow:} \texttt{QR-A-T4.1}/\texttt{QR-B-L3.2}
\(\to\) T-DOME II internal algebra.\par
\noindent\textbf{Source Claim ID:} \texttt{QR-A-T4.1},
\texttt{QR-B-L3.2}.\par
\noindent\textbf{Source current status:} Both selection/necessity
theorems are \texttt{RETRACT}; the revised papers retain Clifford
representations as model choices.\par
\noindent\textbf{Target usage description:} T-DOME II explicitly adopts
\(Cl(V,q)\) and a representation map as supplied model inputs.\par
\noindent\textbf{Arrow status:} \textbf{MODEL}.\par
\noindent\textbf{Failure effect if mis-graded:} A chosen representation
would be advertised as a derived necessity for persistent agents.\par
\noindent\textbf{Notes:} Local Clifford examples exist, but there is no
end-to-end read-out/control/archive common model and no theorem-level
inheritance.\par

\subsubsection*{E5. Two posited Clifford routes to QR-C convergence}
\noindent\textbf{Arrow:} \texttt{QR-A-T4.1}/\texttt{QR-B-L3.2}
\(\to\) Q-RAIF C convergence.\par
\noindent\textbf{Source Claim ID:} \texttt{QR-A-T4.1},
\texttt{QR-B-L3.2}.\par
\noindent\textbf{Source current status:} Both are \texttt{RETRACT} as
selection theorems.\par
\noindent\textbf{Target usage description:} Formerly claimed independent
convergence on Clifford structure.\par
\noindent\textbf{Arrow status:} \textbf{INVALID-OPEN}.\par
\noindent\textbf{Failure effect if mis-graded:} Two shared modelling
choices would masquerade as independent physical convergence.\par
\noindent\textbf{Notes:} Q-RAIF C now proves only an image statement
after a homomorphism is supplied.\par
\noindent\textbf{Repair path:} Independently derive both source
representations and construct a physical map preserving the required
quadratic/signature data.\par

\subsubsection*{E6. T-DOME I claims to T-DOME II motivation}
\noindent\textbf{Arrow:} \texttt{TD-I-ceiling}/\texttt{TD-I-adv}
\(\to\) T-DOME II motivation.\par
\noindent\textbf{Source Claim ID:} \texttt{TD-I-ceiling},
\texttt{TD-I-adv}.\par
\noindent\textbf{Source current status:} The ceiling is
\texttt{RETRACT}; non-Markovian advantage is \texttt{DOWNGRADE} and
depends on initial correlations.\par
\noindent\textbf{Target usage description:} Motivates an agent attempting
to exploit correlation-aware memory.\par
\noindent\textbf{Arrow status:} \textbf{HEUR}.\par
\noindent\textbf{Failure effect if mis-graded:} Memory would again be
called universally necessary for persistence.\par
\noindent\textbf{Notes:} No single protocol proves the claimed
necessity across the two papers.\par

\subsubsection*{E7. T-DOME II necessity to later drift/foreground uses}
\noindent\textbf{Arrow:} \texttt{TD-II-T3.5}
\(\to\) T-DOME III drift and Foundation I foreground narrative.\par
\noindent\textbf{Source Claim ID:} \texttt{TD-II-T3.5}.\par
\noindent\textbf{Source current status:} \texttt{THM/RETRACT}; the
dimension count and optimization/SSB necessity are unsupported.\par
\noindent\textbf{Target usage description:} Supplies only a frame-model
and foreground/background motivation.\par
\noindent\textbf{Arrow status:} \textbf{HEUR}.\par
\noindent\textbf{Failure effect if mis-graded:} Later papers would
inherit a false symmetry-breaking necessity theorem.\par
\noindent\textbf{Notes:} The downstream Gaussian/frame examples stand
only on their own stated assumptions.\par

\subsubsection*{E8. Foundation hull to Consciousness shell}
\noindent\textbf{Arrow:} \texttt{FI-hull}
\(\to\) \texttt{CI-Amb-shell}.\par
\noindent\textbf{Source Claim ID:} \texttt{FI-hull}.\par
\noindent\textbf{Source current status:} \texttt{THM/OK-cond} under a
nonempty invariant class plus modular-invariance and expectation
hypotheses.\par
\noindent\textbf{Target usage description:} Consciousness I supplies
shell generation, closure, nesting, core, and standardness as a new
cross-paper input.\par
\noindent\textbf{Arrow status:} \textbf{IAX}.\par
\noindent\textbf{Failure effect if mis-graded:} Shell covariance and
scaling would be falsely inherited from Foundation I rather than derived
from the independent dynamical axiom \((C)\).\par
\noindent\textbf{Notes:} The interface excludes covariance; no
nontrivial bounded-generator finite-dimensional common model exists.\par

\subsubsection*{E9. Foundation hulls to Brain workspace}
\noindent\textbf{Arrow:} \texttt{FI-hull}/\texttt{FII-share}
\(\to\) Brain Anchor.\par
\noindent\textbf{Source Claim ID:} \texttt{FI-hull},
\texttt{FII-share}.\par
\noindent\textbf{Source current status:} \texttt{FI-hull} is
\texttt{OK-cond}; the former one-state \texttt{FII-share} theorem is
\texttt{RETRACT} and was rebuilt with independent source states.\par
\noindent\textbf{Target usage description:} Provides analogy for a
distinct internal workspace \(\mathcal W_{\mathrm{int}}\).\par
\noindent\textbf{Arrow status:} \textbf{HEUR}.\par
\noindent\textbf{Failure effect if mis-graded:} Source hull
\(\mathcal H_{\mathrm{st}}\), external sector \(\mathcal H_{12}\), and
cognitive workspace \(\mathcal W_{\mathrm{int}}\) would be silently
identified.\par
\noindent\textbf{Notes:} Brain solves non-vacuity in its own finite
witness and imports no Foundation theorem.\par

\subsubsection*{E10. Final-cut diagnosis to CIII framing}
\noindent\textbf{Arrow:} \texttt{HAFF-G-L3.1}
\(\to\) CIII qualitative incompleteness framing.\par
\noindent\textbf{Source Claim ID:} \texttt{HAFF-G-L3.1}.\par
\noindent\textbf{Source current status:} \texttt{META/RETRACT} as a
formal non-self-grounding lemma; programme-level diagnosis only.\par
\noindent\textbf{Target usage description:} Motivates an open
self-description question.\par
\noindent\textbf{Arrow status:} \textbf{HEUR}.\par
\noindent\textbf{Failure effect if mis-graded:} CIII would again claim to
quantify an unproved Paper~G theorem.\par
\noindent\textbf{Notes:} No fixed-point or observation-model bridge is
available.\par

\subsubsection*{E11. Algebraic modular time to phenomenal interpretation}
\noindent\textbf{Arrow:} \texttt{CI-T4.8}/\texttt{CI-T5.7}
\(\to\) CIII phenomenal interpretation.\par
\noindent\textbf{Source Claim ID:} \texttt{CI-T4.8},
\texttt{CI-T5.7}.\par
\noindent\textbf{Source current status:} \texttt{CI-T4.8} is
\texttt{OK-cond}; \texttt{CI-T5.7} is \texttt{DOWNGRADE} because its
reverse direction is axiom \((SFa)\).\par
\noindent\textbf{Target usage description:} Offers an interpretive
reading of algebraic modular flow as relevant to phenomenal time.\par
\noindent\textbf{Arrow status:} \textbf{HEUR}.\par
\noindent\textbf{Failure effect if mis-graded:} Algebraic modular
parameterization would be equated with phenomenal temporality.\par
\noindent\textbf{Notes:} The operator-algebraic and phenomenal domains
have no common validating model.\par

\subsubsection*{E12. CHM remainder to a CIII coefficient tower}
\noindent\textbf{Arrow:} \texttt{CII-P4.1} and its
\(O(R^{d+3})\) remainder \(\to\) CIII coefficient/Fisher tower.\par
\noindent\textbf{Source Claim ID:} \texttt{CII-P4.1}.\par
\noindent\textbf{Source current status:} \texttt{PROP/OK-cond} for CHM
balls with smooth local stress tensor in the declared linearized regime.\par
\noindent\textbf{Target usage description:} Formerly treated a remainder
order as an infinite sequence of independent self-description
directions.\par
\noindent\textbf{Arrow status:} \textbf{INVALID-OPEN}.\par
\noindent\textbf{Failure effect if mis-graded:} A big-\(O\) remainder
would be converted into a normalized physical Fisher spectrum without a
construction.\par
\noindent\textbf{Notes:} The leading stress-tensor extraction survives;
the tower does not.\par
\noindent\textbf{Repair path:} Construct independent normalized
coefficient perturbations, a renormalized kernel, and a typed Fisher map
from CII data to the CIII statistical model.\par

\subsubsection*{E13. Ego-prior model to CIII van Trees use}
\noindent\textbf{Arrow:} \texttt{TD-III-ego}
\(\to\) CIII van Trees setup.\par
\noindent\textbf{Source Claim ID:} \texttt{TD-III-ego}.\par
\noindent\textbf{Source current status:} \texttt{MODEL/OK-cond} for a
specified regular prior and parameterization.\par
\noindent\textbf{Target usage description:} Supplies a possible prior
model for a future typed estimation problem, not a reconstruction floor.\par
\noindent\textbf{Arrow status:} \textbf{MODEL}.\par
\noindent\textbf{Failure effect if mis-graded:} A supplied prior would be
mistaken for a derived self-reference law or positive reconstruction
loss.\par
\noindent\textbf{Notes:} CIII currently has no common sampling model
connecting this prior to its state-reconstruction loss.\par

\subsubsection*{E14. Retracted CIII drift floor to IC formal tier}
\noindent\textbf{Arrow:} \texttt{CIII-T3.3}
\(\to\) Imprint Chain formal tier.\par
\noindent\textbf{Source Claim ID:} \texttt{CIII-T3.3}.\par
\noindent\textbf{Source current status:} \texttt{THM/RETRACT}; a
lossless-orbit code gives zero reconstruction loss.\par
\noindent\textbf{Target usage description:} Formerly supplied a positive
finite-self-model drift floor as a formal roadmap endpoint.\par
\noindent\textbf{Arrow status:} \textbf{INVALID-OPEN}.\par
\noindent\textbf{Failure effect if mis-graded:} The roadmap would
reintroduce a theorem contradicted by an explicit channel construction.\par
\noindent\textbf{Notes:} The formal tier contains no live drift-floor
claim.\par
\noindent\textbf{Repair path:} Prove a new rate-distortion or
single/few-copy estimation theorem under a typed observation model that
excludes lossless orbit coding.\par

\subsubsection*{E15. Retracted physical control order to horizon endpoint}
\noindent\textbf{Arrow:} \texttt{CIII-kn}
\(\to\) Imprint Chain finite-horizon endpoint.\par
\noindent\textbf{Source Claim ID:} \texttt{CIII-kn}.\par
\noindent\textbf{Source current status:} \texttt{THM/RETRACT}; the
normalization, inverse-Fisher budget, and coefficient tower are not
physically sourced.\par
\noindent\textbf{Target usage description:} Formerly supplied
\(k(n)=O(\log n/\log\log n)\) as a self-description theorem.\par
\noindent\textbf{Arrow status:} \textbf{INVALID-OPEN}.\par
\noindent\textbf{Failure effect if mis-graded:} A conditional arithmetic
lemma would be mislabeled as a physical self-model budget.\par
\noindent\textbf{Notes:} Only \texttt{CIII-asym} survives as an abstract
implication.\par
\noindent\textbf{Repair path:} Independently derive a physical normalized
spectrum, justify the resource budget, and construct the coefficient
tower before reconnecting the endpoint.\par

\subsubsection*{E16. Retracted CII gravity package to programme headline}
\noindent\textbf{Arrow:} \texttt{CII-P4.3}
\(\to\) programme gravity headline.\par
\noindent\textbf{Source Claim ID:} \texttt{CII-P4.3}.\par
\noindent\textbf{Source current status:} \texttt{PROP/RETRACT};
derivative order, transversality, and modular-object typing are
insufficient.\par
\noindent\textbf{Target usage description:} Formerly supported an
unconditional gravitational derivation headline.\par
\noindent\textbf{Arrow status:} \textbf{INVALID-OPEN}.\par
\noindent\textbf{Failure effect if mis-graded:} Conditional coefficient
matching would become a general Einstein-dynamics theorem.\par
\noindent\textbf{Notes:} The restricted \texttt{CII-P4.1}
stress-tensor extraction is a different surviving node.\par
\noindent\textbf{Repair path:} Rebuild the geometric axiom package with
restricted derivative order, transversality, distinct modular objects,
and a non-circular geometric functional construction.\par

\subsubsection*{E17. Imprint heuristic to RG and Stable Hulls}
\noindent\textbf{Arrow:} \texttt{WI-princ}
\(\to\) RG Filtering and Stable Hulls.\par
\noindent\textbf{Source Claim ID:} \texttt{WI-princ}.\par
\noindent\textbf{Source current status:} \texttt{HEUR/DOWNGRADE};
\(\operatorname{Imp}_{\mathcal D}\) is decoder-relative and may be zero.\par
\noindent\textbf{Target usage description:} RG offers a heuristic
filtering rereading; Stable Hulls proposes a conditional cognitive
integration bridge.\par
\noindent\textbf{Arrow status:} \textbf{HEUR}.\par
\noindent\textbf{Failure effect if mis-graded:} The downstream papers
would inherit universal nonzero carrying or a formal theorem from the
imprint heuristic.\par
\noindent\textbf{Notes:} Gate 2 and Gate 3 synchronized both downstream
surfaces; the Stable Hulls conclusion now explicitly allows small or null
decoder-relative carrying.\par

\section{How the Papers Fit Together}

The chain can now be mapped directly onto existing and planned papers.

\begin{center}
\begin{tabular}{P{0.34\linewidth}P{0.56\linewidth}}
\toprule
\textbf{Chain segment} & \textbf{Primary paper location} \\
\midrule
source algebra $\to$ coarse-graining & Genesis Trilogy; Foundation~I \\
coarse-graining $\to$ scale-stable imprints & \emph{Renormalization as Imprint Filtering} \\
scale-stable imprints $\to$ imprint-field & \emph{The World as Imprint} \\
imprint-field $\to\mathcal W_{\mathrm{int}}\to$ self-model & \emph{Stable Hulls as Imprint Integration} \S4--\S6; Brain Anchor as candidate; Foundation~II's \(\mathcal H_{12}\) is a distinct external object \\
self-model $\to$ self-reading threshold & \emph{The World as Imprint} \\
self-reading threshold $\dashrightarrow$ finite self-description question & \emph{Finite Self-Models and the Open Explanatory Gap}: arithmetic lemma only; physical edge retracted/open \\
\bottomrule
\end{tabular}
\end{center}

The key point is that the bridge papers do not create a fourth series. They close local gaps between already existing clusters of work. Their natural place is interstitial rather than foundational in their own right. Papers 1 and 2 are therefore best treated as Zenodo working papers serving as interpretive companions to the present roadmap; they are not offered as independent technical contributions to renormalization theory or to cognitive science.

\section{Reading Paths for External Readers}

The body of work is now large enough that not every reader should be asked to read it all in chronological order. Three paths are therefore useful.

\subsection{Two-paper path}

Readers mainly interested in the philosophical architecture can begin with:

\begin{enumerate}
\item \emph{The World as Imprint}
\item \emph{The Imprint Chain}
\end{enumerate}

This path gives the overall frame and the internal grading without requiring the full technical background.

\subsection{Five-paper path}

Readers interested in structure and mind, but not in all of Genesis, can follow:

\begin{enumerate}
\item \emph{The World as Imprint}
\item \emph{Renormalization as Imprint Filtering}
\item \emph{Stable Hulls as Imprint Integration}
\item \emph{The Brain as a Plastic Observable-Algebra Network: The Conditional Fixed-Point Workspace \(\mathcal W_{\mathrm{int}}\)}
\item \emph{Finite Self-Models and the Open Explanatory Gap}
\end{enumerate}

This path is thematically concentrated but is not self-contained. Before using
Consciousness Paper~III, readers must also consult Consciousness Paper~II for
the small-ball remainder, T-DOME Paper~III for the ego-prior model, HAFF
Paper~G for programme-level motivation, and Consciousness Paper~I for the
\((Amb)\), \((C)\), and \((SFa)\) interfaces.

\subsection{Full path}

Readers pursuing the complete architecture should read:

\begin{enumerate}
\item the Genesis Trilogy
\item the Consciousness Series
\item \emph{The World as Imprint}
\item the two bridge papers
\item the present roadmap
\end{enumerate}

The present paper is not the first paper for such readers. It is the paper that helps them not get lost once the archive becomes large.

\section{What the Chain Does Not Claim}

The map is useful only if it also states its non-claims clearly.

First, the chain is not all theorem-level. It contains formal results, bridges, interfaces, and resonances.

Second, the chain does not prove that physics validates Buddhism or any other contemplative tradition. Such resonances are explicitly reader-facing and non-evidential.

Third, the chain does not solve the hard problem. Even if self-reading architecture is clarified, why experience feels like something remains open.

Fourth, the chain does not establish soul survival, identity transfer, or complete self-transparency.

Fifth, the existence of a map does not mean the map is complete. The point is not closure by proclamation, but intelligibility by segmentation.

\section{Conclusion: Displayed, Not Proved}

The point of this paper has been organizational from the start. A large body of work can fail not because its parts are individually weak, but because readers cannot tell how the parts are supposed to hang together. The imprint chain is an attempt to prevent that failure.

What emerges from the map is a cleaner picture of the whole arc. The
Genesis papers study effective structure under accessibility and
coarse-graining. The first bridge paper gives a heuristic RG rereading.
\emph{The World as Imprint} supplies decoder-relative vocabulary. The
second bridge paper proposes \(\mathcal W_{\mathrm{int}}\) as an
internal integration interface. Later consciousness papers study
self-modeling, while a physical finite-self-description bound remains
open after the Paper~III retractions.

This does not amount to a single theorem. It amounts to a graded route from emergence to self-reading. That is enough. A roadmap should not pretend to be the terrain. Its task is to make the terrain traversable.

\begin{thebibliography}{99}

\bibitem{WorldImprint}
S.~Liu, \emph{The World as Imprint: From Trace-Carrying Structure to Self-Reading Mind}, Zenodo (2026), DOI: 10.5281/zenodo.20122718.

\bibitem{HAFFA}
S.~Liu, \emph{Emergent Geometry from Coarse-Grained Observable Algebras: The Holographic Alaya-Field Framework}, Zenodo (2026), DOI: 10.5281/zenodo.18361706.

\bibitem{HAFFD}
S.~Liu, \emph{Observable-Algebra Selection and Emergent Gravitational Phenomena: A Conditional Structural Analysis}, Zenodo (2026), DOI: 10.5281/zenodo.18388881.

\bibitem{FoundationI}
S.~Liu, \emph{Foundation I: From Source to Manifestation --- A Conditional Same-Root Program for Geometry and Modular Flow}, Zenodo (2026), DOI: 10.5281/zenodo.19627786.

\bibitem{FoundationII}
S.~Liu, \emph{Foundation II: Shared Manifestation --- A Conditional Two-Source Program for Shared Accessible Sectors}, Zenodo (2026), DOI: 10.5281/zenodo.19641372.

\bibitem{FoundationIII}
S.~Liu, \emph{Foundation III: The Perfuming Loop --- A Conditional One-Step Source-Level Feedback Program}, Zenodo (2026), DOI: 10.5281/zenodo.19642189.

\bibitem{BrainAnchor}
S.~Liu, \emph{The Brain as a Plastic Observable-Algebra Network: The Conditional Fixed-Point Workspace $\mathcal{W}_{\mathrm{int}}$}, Zenodo (2026), DOI: 10.5281/zenodo.20044329.

\bibitem{PaperIII}
S.~Liu, \emph{Finite Self-Models and the Open Explanatory Gap: A Conditional Arithmetic Analysis}, Zenodo (2026), DOI: 10.5281/zenodo.19607724.

\bibitem{Ladyman2007}
J.~Ladyman and D.~Ross,
\emph{Every Thing Must Go: Metaphysics Naturalized},
Oxford University Press (2007).

\bibitem{French2014}
S.~French,
\emph{The Structure of the World: Metaphysics and Representation},
Oxford University Press (2014).

\bibitem{Batterman2002}
R.~W.~Batterman,
\emph{The Devil in the Details: Asymptotic Reasoning in Explanation, Reduction, and Emergence},
Oxford University Press (2002).

\end{thebibliography}

\clearpage
\section{Global Arrow Summary}

\begin{center}
\begin{tabular}{P{0.20\linewidth}P{0.20\linewidth}P{0.22\linewidth}P{0.18\linewidth}}
\toprule
\textbf{Status} & \textbf{Narrative arrows} & \textbf{Dependency edges} & \textbf{Combined} \\
\midrule
THM & 0 & 0 & 0 \\
COND & 0 & 0 & 0 \\
MODEL & 0 & 2 & 2 \\
IAX & 0 & 1 & 1 \\
DEF & 0 & 0 & 0 \\
HEUR & 6 & 8 & 14 \\
META & 0 & 0 & 0 \\
INVALID-OPEN & 1 & 6 & 7 \\
\midrule
\textbf{Total} & \textbf{7} & \textbf{17} & \textbf{24} \\
\bottomrule
\end{tabular}
\end{center}

The zero counts for THM and COND do not mean that the programme contains
no theorems. They mean that no cross-paper arrow in this roadmap is itself
a verified or conditional theorem. Formal results remain at their own
nodes and do not propagate automatically across interfaces.

The HAFF--Q-RAIF--T-DOME route is therefore not a theorem chain.
HAFF-G to the Q-RAIF no-view diagnosis is HEUR (E3);
Q-RAIF's Clifford choices to T-DOME II are MODEL (E4);
T-DOME I to II and T-DOME II to later drift/foreground uses are HEUR
(E6--E7). Each target must stand on its own hypotheses, model inputs, and
proofs. No status is inherited transitively from an earlier edge.

The Foundation--Consciousness interfaces are likewise local: FI to the
Consciousness I shell is an explicit IAX (E8), while the Foundation-to-Brain
comparison is HEUR (E9). The imprint narrative is a sequence of six HEUR
arrows ending in one INVALID-OPEN physical-horizon question. Seven
INVALID-OPEN entries (P7, E1, E5, E12, E14--E16) each state a repair path;
none is repaired by rhetoric or by the existence of this map.

The main-chain reliability assessment is consequently
\emph{architecturally useful but inferentially non-transitive}. Readers
may use the roadmap to locate compatible questions and model choices, but
must return to each source paper for its domain, assumptions, and proof.
The five-paper path is explicitly not self-contained.

\end{document}
